% Generated draft: requires mathematical and source review.
\documentclass[11pt]{article}
\usepackage[margin=25mm]{geometry}
\usepackage{fontspec}
\setmainfont{Latin Modern Roman}
\usepackage{amsmath,amssymb,mathtools,unicode-math}
\setmathfont{Latin Modern Math}
\usepackage{enumitem,needspace,xurl,hyperref}
\setlength{\parindent}{0pt}
\setlength{\parskip}{5pt}
\setlength{\emergencystretch}{3em}
\newcommand{\N}{\mathbb{N}}
\newcommand{\Z}{\mathbb{Z}}
\newcommand{\Q}{\mathbb{Q}}
\newcommand{\R}{\mathbb{R}}
\newcommand{\C}{\mathbb{C}}
\newcommand{\F}{\mathbb{F}}
\newcommand{\A}{\mathbb{A}}
\newcommand{\PP}{\mathbb{P}}
\newcommand{\E}{\mathbb{E}}
\newcommand{\eps}{\varepsilon}
\newcommand{\dd}{\,\mathrm{d}}
\begin{document}
\section*{UnsolvedMath ID: 30001847}
\textbf{Source problem code:} catalogue entry
\subsection{Definitions and mathematical statement}
The supplied record does not define the objects $G$ and $M$, their admissible range, or the marking they specify on an elliptic curve. It also leaves the base field, the moduli functor, compactification, and the choice between a moduli stack and a coarse moduli space unspecified. Consequently, the classification question cannot yet be stated as an unconditional, fully specified mathematical problem.

For terminology only, given a field $k$ and positive integers $a_0,\ldots,a_n$, the weighted projective scheme is
\[
\mathbb{P}_k(a_0,\ldots,a_n)=\operatorname{Proj} k[x_0,\ldots,x_n],\qquad \deg x_i=a_i.
\]
A weighted complete intersection over $k$, in the scheme interpretation, is a closed subscheme presented as
\[
\operatorname{Proj}\bigl(k[x_0,\ldots,x_n]/(f_1,\ldots,f_c)\bigr),
\]
where the $f_i$ are weighted homogeneous polynomials forming a regular sequence. These definitions explain standard scheme terminology; they do not resolve whether the source intends schemes or stacks, or requires additional smoothness or other conditions.

\textbf{Conditional formulation of the supplied question.} Once an admissible class $\mathcal{A}$ of pairs $(G,M)$ and the associated moduli objects $X_{G,M}$ of marked elliptic curves have been defined, classify all $(G,M)\in\mathcal{A}$ for which $X_{G,M}$ is isomorphic to a weighted projective space, or admits a realization as a complete intersection in a weighted projective space. Under a scheme interpretation over a fixed field $k$, the first property would mean that there exist $n\geq 0$ and positive integers $a_0,\ldots,a_n$ such that
\[
X_{G,M}\cong\mathbb{P}_k(a_0,\ldots,a_n),
\]
and the second would mean that $X_{G,M}$ has a presentation of the weighted complete intersection form above. The symbols $\mathcal{A}$ and $X_{G,M}$ are placeholders for missing source definitions, not supplied constructions.

The source also asks how far this classification has progressed and what further applications these moduli spaces have. It specifies neither a criterion for measuring classification progress nor a particular application to establish. These parts remain research directions rather than determinate mathematical assertions.
\subsection{Short English statement}
Determine which pairs $(G,M)$ give moduli objects of marked elliptic curves that are weighted projective spaces or weighted complete intersections, and investigate their further applications. The supplied record does not define $(G,M)$, the markings, or the precise moduli objects, so the formulation remains conditional. Its open status and claim that the classification is incomplete are supplied dated assessments, not independently verified facts.
\subsection{Sources}
\begin{itemize}
\item Local UnsolvedMath registry, permanent ID $30001847$, title \emph{Weighted-Projective Moduli of Marked Elliptic Curves}; collection: Oberwolfach Reports Open Problems; exact source problem code: \texttt{catalogue entry}.
\item Supplied report citation: \emph{Explicit Methods in Number Theory (2012)}, Oberwolfach Reports, pp.~1971--2044, DOI \texttt{10.4171/owr/2011/35}. Supplied source URL and literature-source link labelled ``original OWR report'': \url{https://doi.org/10.4171/owr/2011/35}. The background describes the report as documenting workshop discussions and open problems. No individual proposer is identified.
\item The registry gives the proposed year as $2011$, while its report citation names $2012$; this discrepancy is preserved without resolution.
\item Literature assessment checked on \texttt{2026-08-21}: the supplied curation labels the problem open and states that the classification of marked-curve moduli spaces that are weighted projective spaces or weighted complete intersections is incomplete. Its listed checked source is the original OWR report linked above. This is supplied dated curation, not an independently verified literature assessment; the record itself requests independent verification of time-sensitive status and cited claims.
\item The supplied statement status is \texttt{corrected\_verified}. Its accompanying verification note reports that the mathematical content was checked against the official Oberwolfach report and that notation, context, and question form were normalized without changing the source-level claim. This records the supplied verification claim; no independent source consultation is asserted here.
\end{itemize}
\end{document}
